% Options for packages loaded elsewhere
\PassOptionsToPackage{unicode}{hyperref}
\PassOptionsToPackage{hyphens}{url}
\documentclass[
  british,
  a4paper,
]{article}
\usepackage{xcolor}
\usepackage[margin=1in]{geometry}
\usepackage{amsmath,amssymb}
\setcounter{secnumdepth}{-\maxdimen} % remove section numbering
\usepackage{iftex}
\ifPDFTeX
  \usepackage[T1]{fontenc}
  \usepackage[utf8]{inputenc}
  \usepackage{textcomp} % provide euro and other symbols
\else % if luatex or xetex
  \usepackage{unicode-math} % this also loads fontspec
  \defaultfontfeatures{Scale=MatchLowercase}
  \defaultfontfeatures[\rmfamily]{Ligatures=TeX,Scale=1}
\fi
\usepackage{lmodern}
\ifPDFTeX\else
  % xetex/luatex font selection
  \setmainfont[]{Noto Serif}
  \setmonofont[]{Noto Sans Mono}
\fi
% Use upquote if available, for straight quotes in verbatim environments
\IfFileExists{upquote.sty}{\usepackage{upquote}}{}
\IfFileExists{microtype.sty}{% use microtype if available
  \usepackage[]{microtype}
  \UseMicrotypeSet[protrusion]{basicmath} % disable protrusion for tt fonts
}{}
\makeatletter
\@ifundefined{KOMAClassName}{% if non-KOMA class
  \IfFileExists{parskip.sty}{%
    \usepackage{parskip}
  }{% else
    \setlength{\parindent}{0pt}
    \setlength{\parskip}{6pt plus 2pt minus 1pt}}
}{% if KOMA class
  \KOMAoptions{parskip=half}}
\makeatother
\usepackage{longtable,booktabs,array}
\newcounter{none} % for unnumbered tables
\usepackage{calc} % for calculating minipage widths
% Correct order of tables after \paragraph or \subparagraph
\usepackage{etoolbox}
\makeatletter
\patchcmd\longtable{\par}{\if@noskipsec\mbox{}\fi\par}{}{}
\makeatother
% Allow footnotes in longtable head/foot
\IfFileExists{footnotehyper.sty}{\usepackage{footnotehyper}}{\usepackage{footnote}}
\makesavenoteenv{longtable}
% definitions for citeproc citations
\NewDocumentCommand\citeproctext{}{}
\NewDocumentCommand\citeproc{mm}{%
  \begingroup\def\citeproctext{#2}\cite{#1}\endgroup}
\makeatletter
 % allow citations to break across lines
 \let\@cite@ofmt\@firstofone
 % avoid brackets around text for \cite:
 \def\@biblabel#1{}
 \def\@cite#1#2{{#1\if@tempswa , #2\fi}}
\makeatother
\newlength{\cslhangindent}
\setlength{\cslhangindent}{1.5em}
\newlength{\csllabelwidth}
\setlength{\csllabelwidth}{3em}
\newenvironment{CSLReferences}[2] % #1 hanging-indent, #2 entry-spacing
 {\begin{list}{}{%
  \setlength{\itemindent}{0pt}
  \setlength{\leftmargin}{0pt}
  \setlength{\parsep}{0pt}
  % turn on hanging indent if param 1 is 1
  \ifodd #1
   \setlength{\leftmargin}{\cslhangindent}
   \setlength{\itemindent}{-1\cslhangindent}
  \fi
  % set entry spacing
  \setlength{\itemsep}{#2\baselineskip}}}
 {\end{list}}
\usepackage{calc}
\newcommand{\CSLBlock}[1]{\hfill\break\parbox[t]{\linewidth}{\strut\ignorespaces#1\strut}}
\newcommand{\CSLLeftMargin}[1]{\parbox[t]{\csllabelwidth}{\strut#1\strut}}
\newcommand{\CSLRightInline}[1]{\parbox[t]{\linewidth - \csllabelwidth}{\strut#1\strut}}
\newcommand{\CSLIndent}[1]{\hspace{\cslhangindent}#1}
\ifLuaTeX
\usepackage[bidi=basic,shorthands=off]{babel}
\else
\usepackage[bidi=default,shorthands=off]{babel}
\fi
\ifPDFTeX
\else
\babelfont{rm}[]{Noto Serif}
\fi
\ifLuaTeX
  \usepackage{selnolig} % disable illegal ligatures
\fi
\setlength{\emergencystretch}{3em} % prevent overfull lines
\providecommand{\tightlist}{%
  \setlength{\itemsep}{0pt}\setlength{\parskip}{0pt}}
\usepackage{bookmark}
\IfFileExists{xurl.sty}{\usepackage{xurl}}{} % add URL line breaks if available
\urlstyle{same}
\hypersetup{
  pdftitle={Germanic lexeme report assembly pilot 01},
  pdflang={en-GB},
  hidelinks,
  pdfcreator={LaTeX via pandoc}}

\title{Germanic lexeme report assembly pilot 01}
\usepackage{etoolbox}
\makeatletter
\providecommand{\subtitle}[1]{% add subtitle to \maketitle
  \apptocmd{\@title}{\par {\large #1 \par}}{}{}
}
\makeatother
\subtitle{Representative model-entry assembly test}
\author{}
\date{}

\begin{document}
\maketitle

This pilot assembles a small representative subset of the current
Germanic model entries for publication-format testing. It preserves
entry prose, Pandoc citations, and tables from the source
\texttt{.model.md} files while normalizing heading levels only in this
assembled copy.

The included entries are listed in \texttt{pilot\_manifest.tsv} and
appear here in stable manifest order. Implementation reports, reviewer
checklists, source ledgers, packets, and research notes are
intentionally excluded from the assembled body.

\subsection{bake --- OE bacan}\label{bake-oe-bacan}

{\def\LTcaptype{none} % do not increment counter
\begin{longtable}[]{@{}ll@{}}
\toprule\noalign{}
Field & Value \\
\midrule\noalign{}
\endhead
\bottomrule\noalign{}
\endlastfoot
PROTO & *bákaną \\
PROTOFORM & *bákaną \\
COUNTERPART & bacan \\
DERIVATION\_CLASS & regular \\
\end{longtable}
}

\subsubsection{Transducer input and
output}\label{transducer-input-and-output}

{\def\LTcaptype{none} % do not increment counter
\begin{longtable}[]{@{}ll@{}}
\toprule\noalign{}
Item & Value \\
\midrule\noalign{}
\endhead
\bottomrule\noalign{}
\endlastfoot
lexical item & bake \\
citation reconstruction / lexeme label & *bákaną \\
selected input form & *bákaną \\
Old English target & bacan \\
classification & regular \\
documented output & *bákaną -\textgreater{} bacan \\
\end{longtable}
}

\subsubsection{Reconstruction and comparative
evidence}\label{reconstruction-and-comparative-evidence}

Orel reconstructs the verb as \texttt{*bakanan} and cites Old English
\texttt{bacan} beside Old High German \texttt{backan,\ bahhan} (Orel
2003). Campbell gives \texttt{bacan} as one of the standard examples of
Old English A-restoration before a single consonant, and Ringe and
Taylor state the same development from \texttt{*bakan} to Old English
\texttt{bacan} (Campbell 1959, 61; Ringe and Taylor 2014).

\subsubsection{Old English evidence}\label{old-english-evidence}

Bosworth-Toller and Clark Hall both record \texttt{bacan} as the
ordinary Old English verb `to bake' (Bosworth and Toller 1898, 72; Clark
Hall 1960). The target in this entry is therefore the attested
infinitive headword itself, not a selected oblique or finite paradigm
cell.

\subsubsection{Development to Old
English}\label{development-to-old-english}

From \texttt{*bákaną}, Anglo-Frisian brightening first gives
\texttt{*bækaną}. A-restoration then returns the stem vowel to
\texttt{a} before single \texttt{k} plus the back-vocalic infinitive
suffix, and later apocope and weak-tail reduction yield \texttt{bacan}
(Campbell 1959, 61; Ringe and Taylor 2014). The development is therefore
straightforward: \texttt{*bákaną\ \textgreater{}\ bacan}.

\subsection{bow --- OE bēag}\label{bow-oe-bux113ag}

{\def\LTcaptype{none} % do not increment counter
\begin{longtable}[]{@{}ll@{}}
\toprule\noalign{}
Field & Value \\
\midrule\noalign{}
\endhead
\bottomrule\noalign{}
\endlastfoot
PROTO & *béuganą \\
PROTOFORM & *báug \\
COUNTERPART & bēag \\
DERIVATION\_CLASS & late\_analogy \\
\end{longtable}
}

\subsubsection{Transducer input and
output}\label{transducer-input-and-output-1}

{\def\LTcaptype{none} % do not increment counter
\begin{longtable}[]{@{}ll@{}}
\toprule\noalign{}
Item & Value \\
\midrule\noalign{}
\endhead
\bottomrule\noalign{}
\endlastfoot
lexical item & bow \\
citation reconstruction / lexeme label & *béuganą \\
selected input form & *báug \\
Old English target & bēag \\
classification & late\_analogy \\
documented output & *báug -\textgreater{} bēag \\
\end{longtable}
}

\subsubsection{Reconstruction and comparative
evidence}\label{reconstruction-and-comparative-evidence-1}

The inherited verb belongs to the class-II strong-verb family
\texttt{*béuganą} (Ringe and Taylor 2014, 55). Within that paradigm,
however, the infinitive and the singular preterite continue different
ablaut grades. The selected input \texttt{*báug} is the singular
preterite cell, whereas the citation form \texttt{*béuganą} is the
infinitive.

Campbell's account of Old English class-II strong verbs treats the
singular preterite \texttt{au\ \textgreater{}\ ēa} development as
regular in this environment (Campbell 1959, 53). That is the
phonological path relevant for \texttt{bēag}, whereas the analogical
\texttt{ū} of the present stem belongs to the separate history behind
\texttt{būgan} (Ringe and Taylor 2014, 55).

\subsubsection{Old English evidence}\label{old-english-evidence-1}

Bosworth-Toller and Clark Hall both record \texttt{bēag} as a preterite
form of \texttt{būgan} (Bosworth and Toller 1898, 122; Clark Hall 1960,
45). The form discussed here is therefore an attested Old English verbal
form, not a reconstructed substitute for the infinitive.

The ordinary dictionary headword remains \texttt{būgan}, but the
relevant comparison form for this entry is the singular preterite
\textbf{\texttt{bēag}}. That is the paradigm cell in which the inherited
\texttt{*au} grade is preserved most directly.

\subsubsection{Development to Old
English}\label{development-to-old-english-1}

From \texttt{*báug}, Anglo-Frisian fronting and the later leveling of
the diphthong produce \texttt{bēag} (Campbell 1959, 53). No special
analogical repair is needed for that cell. The form is the regular Old
English outcome of the singular-preterite grade.

The analogical element in the wider lexeme belongs instead to the
present stem seen in \texttt{būgan}. The selected input differs from the
citation form because the regular inherited pathway survives more
transparently in the preterite than in the infinitive.

\subsubsection{Paradigm comparison}\label{paradigm-comparison}

The comparison below is manual. It distinguishes the regular singular
preterite from the more familiar infinitival citation form.

{\def\LTcaptype{none} % do not increment counter
\begin{longtable}[]{@{}
  >{\raggedright\arraybackslash}p{(\linewidth - 8\tabcolsep) * \real{0.2000}}
  >{\raggedright\arraybackslash}p{(\linewidth - 8\tabcolsep) * \real{0.2000}}
  >{\raggedright\arraybackslash}p{(\linewidth - 8\tabcolsep) * \real{0.2000}}
  >{\raggedright\arraybackslash}p{(\linewidth - 8\tabcolsep) * \real{0.2000}}
  >{\raggedright\arraybackslash}p{(\linewidth - 8\tabcolsep) * \real{0.2000}}@{}}
\toprule\noalign{}
\begin{minipage}[b]{\linewidth}\raggedright
PGmc cell / interpretation
\end{minipage} & \begin{minipage}[b]{\linewidth}\raggedright
Candidate input
\end{minipage} & \begin{minipage}[b]{\linewidth}\raggedright
Expected or documented OE outcome
\end{minipage} & \begin{minipage}[b]{\linewidth}\raggedright
OE comparison form
\end{minipage} & \begin{minipage}[b]{\linewidth}\raggedright
Result
\end{minipage} \\
\midrule\noalign{}
\endhead
\bottomrule\noalign{}
\endlastfoot
citation infinitive & *béuganą & inherited present-stem history behind
\texttt{būgan} & būgan & establishes the lexeme, but not the selected
target \\
singular preterite & *báug & compact-trace output: \texttt{bēag} & bēag
& exact match between input, output, and attested cell \\
past participial branch & participial \texttt{*bugan-} type & later
participial outcomes & bogen-type evidence & relevant to the paradigm,
but not the clearest match for this entry \\
\end{longtable}
}

The singular preterite is the relevant comparison form. It gives a
direct lautgesetzlich path to attested \texttt{bēag}, while the citation
form \texttt{būgan} belongs to a paradigm whose present stem has already
undergone later leveling.

\subsection{craft --- OE cræft}\label{craft-oe-cruxe6ft}

{\def\LTcaptype{none} % do not increment counter
\begin{longtable}[]{@{}ll@{}}
\toprule\noalign{}
Field & Value \\
\midrule\noalign{}
\endhead
\bottomrule\noalign{}
\endlastfoot
PROTO & *kráftiz \\
PROTOFORM & *kráftaz \\
COUNTERPART & cræft \\
DERIVATION\_CLASS & early\_analogy \\
\end{longtable}
}

\subsubsection{Transducer input and
output}\label{transducer-input-and-output-2}

{\def\LTcaptype{none} % do not increment counter
\begin{longtable}[]{@{}ll@{}}
\toprule\noalign{}
Item & Value \\
\midrule\noalign{}
\endhead
\bottomrule\noalign{}
\endlastfoot
lexical item & craft \\
citation reconstruction / lexeme label & *kráftiz \\
selected input form & *kráftaz \\
Old English target & cræft \\
classification & early\_analogy \\
documented output & \texttt{*kráftaz\ -\textgreater{}\ cræft} \\
\end{longtable}
}

\subsubsection{Reconstruction and comparative
evidence}\label{reconstruction-and-comparative-evidence-2}

Comparative sources disagree about the older stem class. Kroonen gives a
u-stem \texttt{*kraftu-}, while Orel prints \texttt{*kraftiz\ *kraftuz}
(Kroonen 2013, 340; Orel 2003, 259). The comparative label
\texttt{*kráftiz} remains in view as a lexeme-level shorthand, while
\texttt{*kráftaz} is the pre-Old-English form used for the Old English
derivation.

\subsubsection{Old English evidence}\label{old-english-evidence-2}

The Old English noun itself is secure. Clark Hall and Bosworth-Toller
both give \texttt{cræft} as the headword (Clark Hall 1960, 19; Bosworth
and Toller 1898, 145).

\subsubsection{Development to Old
English}\label{development-to-old-english-2}

The comparison is between possible pre-Old-English inputs. The i-stem
comparator \texttt{*kráftiz} gives \texttt{creft}, while the u-stem
comparator \texttt{*kráftuz} gives \texttt{craft}. The a-stem-shaped
input \texttt{*kráftaz} yields \texttt{cræft} and is therefore the form
used for the Old English derivation. This does not require treating the
comparative dictionaries as identical; it shows the narrower point that
the Old English derivation needs a pre-Old-English form without the
i-umlaut trigger of \texttt{*-iz} and without the back-vowel outcome
associated with the u-stem comparator.

\subsubsection{Form comparison}\label{form-comparison}

{\def\LTcaptype{none} % do not increment counter
\begin{longtable}[]{@{}lll@{}}
\toprule\noalign{}
Candidate input & OE output & Result \\
\midrule\noalign{}
\endhead
\bottomrule\noalign{}
\endlastfoot
*kráftiz & \texttt{creft} & non-match; i-stem comparator \\
*kráftuz & \texttt{craft} & non-match; u-stem comparator \\
*kráftaz & \texttt{cræft} & exact match; selected pre-OE input \\
\end{longtable}
}

\subsection{fowl --- OE fugol}\label{fowl-oe-fugol}

{\def\LTcaptype{none} % do not increment counter
\begin{longtable}[]{@{}ll@{}}
\toprule\noalign{}
Field & Value \\
\midrule\noalign{}
\endhead
\bottomrule\noalign{}
\endlastfoot
PROTO & *fúglaz \\
PROTOFORM & *fúglaz \\
COUNTERPART & fugol \\
DERIVATION\_CLASS & unexplained\_unmodelled \\
\end{longtable}
}

\subsubsection{Transducer input and
output}\label{transducer-input-and-output-3}

{\def\LTcaptype{none} % do not increment counter
\begin{longtable}[]{@{}ll@{}}
\toprule\noalign{}
Item & Value \\
\midrule\noalign{}
\endhead
\bottomrule\noalign{}
\endlastfoot
lexical item & fowl \\
citation reconstruction / lexeme label & *fúglaz \\
selected input form & *fúglaz \\
Old English target & fugol \\
classification & unexplained\_unmodelled \\
computed regular output & *fúglaz -\textgreater{} fogol \\
\end{longtable}
}

\subsubsection{Reconstruction and comparative
evidence}\label{reconstruction-and-comparative-evidence-3}

The noun is the ordinary Germanic a-stem \texttt{*fúglaz}, continued by
forms such as Old Norse \texttt{fugl} and Old High German \texttt{fogal}
(Kroonen 2013, 197; Orel 2003, 155). There is no stem-class or
paradigm-cell dispute behind this entry. The comparative headword and
the selected input are the same.

The difficulty lies instead in the stressed root vowel. Under the
regular West Germanic and Old English development, that \texttt{u}
should lower before the following non-high vowel, yielding an
\texttt{o}-vocalism (Ringe and Taylor 2014, 42--43; Campbell 1959, 43).

\subsubsection{Old English evidence}\label{old-english-evidence-3}

Old English dictionaries record the noun as \texttt{fugol}, with variant
spelling \texttt{fugel} (Bosworth and Toller 1898, 282; Clark Hall 1960,
138). The target is therefore an attested ordinary Old English noun, not
a reconstructed or selectively chosen paradigm form.

The attested word already contains the crucial problem. Its medial
\texttt{-o-} is the ordinary parasite vowel of Old English cluster
phonology, but the root \texttt{fu-} retains \texttt{u} where the
regular history predicts \texttt{fo-} (Campbell 1959, 150).

\subsubsection{Development to Old
English}\label{development-to-old-english-3}

From \texttt{*fúglaz}, the regular cascade yields \texttt{fogol}: the
root vowel lowers before the following non-high vowel (Ringe and Taylor
2014, 42--43; Campbell 1959, 43), final \texttt{z} is lost, and the
cluster is resolved by the usual medial vowel (Ringe and Taylor 2014,
345; Campbell 1959, 150). That is the expected inherited outcome.

The attested Old English noun is \texttt{fugol}, not \texttt{fogol}.
Luick and later handbooks treat this preservation of \texttt{u} as a
small inherited residue, not as a categorical sound law (Luick
1914-\/-1940, 148; Ringe and Taylor 2014, 47). The item therefore
remains a genuine lexical exception rather than the output of a
recoverable regular mechanism.

\subsubsection{Expected and attested
forms}\label{expected-and-attested-forms}

The comparison below is manual. It distinguishes the regular prediction
from the attested Old English noun.

{\def\LTcaptype{none} % do not increment counter
\begin{longtable}[]{@{}
  >{\raggedright\arraybackslash}p{(\linewidth - 4\tabcolsep) * \real{0.3333}}
  >{\raggedright\arraybackslash}p{(\linewidth - 4\tabcolsep) * \real{0.3333}}
  >{\raggedright\arraybackslash}p{(\linewidth - 4\tabcolsep) * \real{0.3333}}@{}}
\toprule\noalign{}
\begin{minipage}[b]{\linewidth}\raggedright
Form
\end{minipage} & \begin{minipage}[b]{\linewidth}\raggedright
Status
\end{minipage} & \begin{minipage}[b]{\linewidth}\raggedright
Relevance to this entry
\end{minipage} \\
\midrule\noalign{}
\endhead
\bottomrule\noalign{}
\endlastfoot
\texttt{fogol} & computed regular output from \texttt{*fúglaz} &
establishes the expected inherited outcome \\
\texttt{fugol} & attested Old English form & selected target; preserves
unexplained root \texttt{u} \\
\texttt{fugel} & attested variant spelling & secondary spelling variant
of the attested noun \\
\end{longtable}
}

The unresolved point lies only in the root vowel. The medial
\texttt{-o-} is regular, but no accepted lautgesetzlich pathway has been
found from \texttt{*fúglaz} to attested \texttt{fugol}.

\subsection{thistle --- OE þistles}\label{thistle-oe-uxfeistles}

{\def\LTcaptype{none} % do not increment counter
\begin{longtable}[]{@{}ll@{}}
\toprule\noalign{}
Field & Value \\
\midrule\noalign{}
\endhead
\bottomrule\noalign{}
\endlastfoot
PROTO & *θéstilaz \\
PROTOFORM & *θístilas \\
COUNTERPART & þistles \\
DERIVATION\_CLASS & late\_analogy \\
\end{longtable}
}

\subsubsection{Transducer input and
output}\label{transducer-input-and-output-4}

{\def\LTcaptype{none} % do not increment counter
\begin{longtable}[]{@{}ll@{}}
\toprule\noalign{}
Item & Value \\
\midrule\noalign{}
\endhead
\bottomrule\noalign{}
\endlastfoot
lexical item & thistle \\
citation reconstruction / lexeme label & *θéstilaz \\
selected input form & *θístilas \\
Old English target & þistles \\
classification & late\_analogy \\
documented output & \texttt{*θístilas\ -\textgreater{}\ þistles} \\
\end{longtable}
}

\subsubsection{Reconstruction and comparative
evidence}\label{reconstruction-and-comparative-evidence-4}

The comparative tradition is divided. Orel prints \texttt{*þe(x)stilaz},
while Kluge-Seebold gives \texttt{*þistila-} (Orel 2003, 458; Kluge
2011). The comparative label \texttt{*θéstilaz} therefore remains in
view as the lexeme-level headword, while the selected input
\texttt{*θístilas} is a specific genitive singular cell.

\subsubsection{Old English evidence}\label{old-english-evidence-4}

The ordinary simplex headword tradition is broken \texttt{þistel} /
\texttt{ðistel}. Clark Hall gives \texttt{ðistel} as the noun headword
(Clark Hall 1960, 326). The selected target here is the genitive
singular \texttt{þistles}, which preserves the same stem in an oblique
form where the cluster is medial.

\subsubsection{Development to Old
English}\label{development-to-old-english-4}

Campbell's discussion of cluster nouns shows the contrast clearly.
Simplex forms often develop a parasite vowel in word-final obstruent +
sonorant clusters, while comparable medial clusters remain unbroken; his
examples include \texttt{hrefn}, \texttt{tacn}, \texttt{wépn}, and
\texttt{botm} beside forms with parasitic vowels elsewhere in the same
lexical class (Campbell 1959, 151). The selected genitive singular
\texttt{*θístilas} therefore supplies the conservative comparison form:
the cluster is medial and the regular development yields
\texttt{þistles}, while the simplex nominative belongs to the broken
headword tradition \texttt{þistel}.

\subsubsection{Paradigm comparison}\label{paradigm-comparison-1}

The comparison below is manual. It shows the contrast between the
citation form and the selected genitive singular cell.

{\def\LTcaptype{none} % do not increment counter
\begin{longtable}[]{@{}
  >{\raggedright\arraybackslash}p{(\linewidth - 8\tabcolsep) * \real{0.2000}}
  >{\raggedright\arraybackslash}p{(\linewidth - 8\tabcolsep) * \real{0.2000}}
  >{\raggedright\arraybackslash}p{(\linewidth - 8\tabcolsep) * \real{0.2000}}
  >{\raggedright\arraybackslash}p{(\linewidth - 8\tabcolsep) * \real{0.2000}}
  >{\raggedright\arraybackslash}p{(\linewidth - 8\tabcolsep) * \real{0.2000}}@{}}
\toprule\noalign{}
\begin{minipage}[b]{\linewidth}\raggedright
PGmc cell / interpretation
\end{minipage} & \begin{minipage}[b]{\linewidth}\raggedright
Candidate input
\end{minipage} & \begin{minipage}[b]{\linewidth}\raggedright
OE output or comparison
\end{minipage} & \begin{minipage}[b]{\linewidth}\raggedright
OE comparison form
\end{minipage} & \begin{minipage}[b]{\linewidth}\raggedright
Result
\end{minipage} \\
\midrule\noalign{}
\endhead
\bottomrule\noalign{}
\endlastfoot
citation nominative singular & *θéstilaz & computed output:
\texttt{þistl} & þistel & useful citation-form background, but not the
selected target \\
selected genitive singular & *θístilas & computed output:
\texttt{þistles} & þistles & exact match for the chosen conservative
cell \\
\end{longtable}
}

\subsection{weapon --- OE wǣpn}\label{weapon-oe-wux1e3pn}

{\def\LTcaptype{none} % do not increment counter
\begin{longtable}[]{@{}ll@{}}
\toprule\noalign{}
Field & Value \\
\midrule\noalign{}
\endhead
\bottomrule\noalign{}
\endlastfoot
PROTO & *wḗpną \\
PROTOFORM & *wḗpną \\
COUNTERPART & wǣpn \\
DERIVATION\_CLASS & regular \\
\end{longtable}
}

\subsubsection{Transducer input and
output}\label{transducer-input-and-output-5}

{\def\LTcaptype{none} % do not increment counter
\begin{longtable}[]{@{}ll@{}}
\toprule\noalign{}
Item & Value \\
\midrule\noalign{}
\endhead
\bottomrule\noalign{}
\endlastfoot
lexical item & weapon \\
citation reconstruction / lexeme label & \texttt{*wḗpną} \\
selected input form & \texttt{*wḗpną} \\
Old English target & \texttt{wǣpn} \\
classification & regular \\
documented output & \texttt{*wḗpną\ -\textgreater{}\ wǣpn} \\
\end{longtable}
}

\subsubsection{Reconstruction and comparative
evidence}\label{reconstruction-and-comparative-evidence-5}

Kroonen reconstructs a double-stem noun
\texttt{*wēbna-\ \textasciitilde{}\ *wēpna-} and cites OE \texttt{wæpn}
among its reflexes (Kroonen 2013, 617). The selected input
\texttt{*wḗpną} represents the unbroken citation-form noun rather than
the later broken simplex.

\subsubsection{Old English evidence}\label{old-english-evidence-5}

Campbell's cluster-noun discussion preserves unbroken \texttt{wépn}
beside broken \texttt{wépen}-type forms (Campbell 1959, 150; 1959,
226--27). Bright contrasts broken nominative \texttt{wǣpen/wapen} with
unbroken oblique \texttt{wǣpnes}, while Clark Hall lemmatizes the noun
under \texttt{wapen} and also preserves unbroken forms in compounds and
related spellings (Bright 1971, 29; Clark Hall 1960, 355).

\subsubsection{Development to Old
English}\label{development-to-old-english-5}

Northwest Germanic lowering gives \texttt{wǣpn}, and loss of the final
nasal vowel leaves the unbroken cluster word-finally. The selected
target is the attested unbroken form \texttt{wǣpn}.

\subsubsection{Form note}\label{form-note}

The ordinary late West Saxon simplex headword is \texttt{wǣpen}, but
Campbell's noun-class discussion also preserves unbroken \texttt{wépn}
beside broken \texttt{wépen}-type forms (Campbell 1959, 150; 1959,
226--27). \texttt{wǣpnes} remains the regular unbroken oblique
comparator (Bright 1971, 29; Clark Hall 1960, 355).

\subsection{will --- OE willa}\label{will-oe-willa}

{\def\LTcaptype{none} % do not increment counter
\begin{longtable}[]{@{}ll@{}}
\toprule\noalign{}
Field & Value \\
\midrule\noalign{}
\endhead
\bottomrule\noalign{}
\endlastfoot
PROTO & *wéljô \\
PROTOFORM & *wéljô \\
COUNTERPART & willa \\
DERIVATION\_CLASS & regular \\
\end{longtable}
}

\subsubsection{Transducer input and
output}\label{transducer-input-and-output-6}

{\def\LTcaptype{none} % do not increment counter
\begin{longtable}[]{@{}ll@{}}
\toprule\noalign{}
Item & Value \\
\midrule\noalign{}
\endhead
\bottomrule\noalign{}
\endlastfoot
lexical item & will \\
citation reconstruction / lexeme label & \texttt{*wéljô} \\
selected input form & \texttt{*wéljô} \\
Old English target & \texttt{willa} \\
classification & regular \\
documented output & \texttt{*wéljô\ -\textgreater{}\ willa} \\
\end{longtable}
}

\subsubsection{Reconstruction and comparative
evidence}\label{reconstruction-and-comparative-evidence-6}

Kroonen separates noun \texttt{*weljan-\ 2} `will, wish' from verb
\texttt{*weljan-\ 1} `to want', while Orel and Kluge represent the noun
as \texttt{*weljōn/*weljOn} (Kroonen 2013; Orel 2003; Kluge 2011). The
selected derivational form \texttt{*wéljô} is the noun-side input used
for this row.

\subsubsection{Old English evidence}\label{old-english-evidence-6}

Clark Hall lemmatizes noun \texttt{willa\ m.} separately from verb
\texttt{willan} (Clark Hall 1960, 368). The selected target is the noun
citation form, not the related verb.

\subsubsection{Development to Old
English}\label{development-to-old-english-6}

From \texttt{*wéljô}, j-gemination yields a heavy stem, i-umlaut gives
\texttt{will-}, and later shortening plus j-loss produce \texttt{willa}.
The noun is therefore a regular weak masculine outcome.

\subsubsection{Lexical note}\label{lexical-note}

The target here is the noun \texttt{willa} `will, wish'. Related verb
\texttt{willan} belongs to a separate lexeme and should not be
substituted for the noun row (Kroonen 2013; Clark Hall 1960, 368).

\subsection{youth --- OE ġeoguþ}\label{youth-oe-ux121eoguuxfe}

{\def\LTcaptype{none} % do not increment counter
\begin{longtable}[]{@{}ll@{}}
\toprule\noalign{}
Field & Value \\
\midrule\noalign{}
\endhead
\bottomrule\noalign{}
\endlastfoot
PROTO & *júgunθiz \\
PROTOFORM & *júgunθ \\
COUNTERPART & ġeoguþ \\
DERIVATION\_CLASS & early\_analogy \\
\end{longtable}
}

\subsubsection{Transducer input and
output}\label{transducer-input-and-output-7}

{\def\LTcaptype{none} % do not increment counter
\begin{longtable}[]{@{}ll@{}}
\toprule\noalign{}
Item & Value \\
\midrule\noalign{}
\endhead
\bottomrule\noalign{}
\endlastfoot
lexical item & youth \\
citation reconstruction / lexeme label & *júgunθiz \\
selected input form & *júgunθ \\
Old English target & ġeoguþ \\
classification & early\_analogy \\
documented output & *júgunθ -\textgreater{} ġeoguþ \\
\end{longtable}
}

\subsubsection{Reconstruction and comparative
evidence}\label{reconstruction-and-comparative-evidence-7}

The wider etymological tradition reconstructs an earlier form of the
word as \texttt{*ju(w)unþi-} (Kroonen 2013; Fulk 2018). The selected
comparative label \texttt{*júgunθiz} already stands at a later Germanic
stage with \texttt{g}, and the chosen input \texttt{*júgunθ} is later
again: it represents the form after final \texttt{-i} has been lost.

That staging matters because Ringe and Taylor explicitly give the
sequence
\texttt{*jugunþi\ \textgreater{}\ *juguþ\ \textgreater{}\ OE\ geoguþ\ \textasciitilde{}\ iuguþ}
(Ringe and Taylor 2014). Campbell makes the same point with the parallel
\texttt{duguþ\ \textless{}\ *dugunþ-}, adding \texttt{geoguþ} to the
same history (Campbell 1959). The selected input therefore differs from
the broader comparative headword because the Old English development
must begin after early loss of final \texttt{-i}.

\subsubsection{Old English evidence}\label{old-english-evidence-7}

The Old English noun is attested with varying spellings. Campbell
records forms of the \texttt{iuguþ\ /\ gioguð\ /\ geoguð} type, and
Ringe and Taylor likewise cite \texttt{geoguþ\ \textasciitilde{}\ iuguþ}
(Campbell 1959; Ringe and Taylor 2014). The form is normalized here as
\texttt{ġeoguþ}: the initial palatal is written with \texttt{ġ}, and the
attested spelling variation is treated as orthographic rather than
lexical.

Nothing in the source stack suggests that a different paradigm cell
should be chosen. The relevant Old English comparison form is the noun
\texttt{ġeoguþ} itself.

\subsubsection{Development to Old
English}\label{development-to-old-english-7}

The decisive early step is the loss of final \texttt{-i} before the Old
English umlaut stage. If that high vowel remained, the word would
develop an over-umlauted \texttt{y}-type vowel instead of the attested
form (Ringe and Taylor 2014).

From the selected input \texttt{*júgunθ}, the later development is
regular: palatal fronting yields \texttt{*jéugunθ}; nasal-spirant
lengthening and loss give \texttt{*jéogūθ}; unstressed long-vowel
shortening then produces \texttt{*jéoguθ}, which surfaces as
\texttt{ġeoguþ}. Medial \texttt{u} remains preserved because the
handbooks treat this environment as one of stem-\texttt{u} harmony after
stressed \texttt{u}, citing forms such as \texttt{munuc},
\texttt{duguþ}, and \texttt{iuguþ} (Campbell 1959; Sievers 1965; Luick
1914-\/-1940, 397).

\subsubsection{Stage comparison}\label{stage-comparison}

The comparison below is manual. It separates the broader comparative
headword from the later stages relevant to the Old English noun.

{\def\LTcaptype{none} % do not increment counter
\begin{longtable}[]{@{}
  >{\raggedright\arraybackslash}p{(\linewidth - 6\tabcolsep) * \real{0.2500}}
  >{\raggedright\arraybackslash}p{(\linewidth - 6\tabcolsep) * \real{0.2500}}
  >{\raggedright\arraybackslash}p{(\linewidth - 6\tabcolsep) * \real{0.2500}}
  >{\raggedright\arraybackslash}p{(\linewidth - 6\tabcolsep) * \real{0.2500}}@{}}
\toprule\noalign{}
\begin{minipage}[b]{\linewidth}\raggedright
Stage / interpretation
\end{minipage} & \begin{minipage}[b]{\linewidth}\raggedright
Candidate form
\end{minipage} & \begin{minipage}[b]{\linewidth}\raggedright
Old English outcome or comparison
\end{minipage} & \begin{minipage}[b]{\linewidth}\raggedright
Relevance to this entry
\end{minipage} \\
\midrule\noalign{}
\endhead
\bottomrule\noalign{}
\endlastfoot
earlier etymological headword & \texttt{*ju(w)unþi-} & comparative
family background & older comparative reconstruction of the lexeme \\
later g-bearing comparative label & \texttt{*júgunθiz} & citation
reconstruction / lexeme label & preserves the later Germanic stage
behind the selected entry \\
selected OE-facing input & \texttt{*júgunθ} & compact-trace output:
\texttt{ġeoguþ} & exact match for the chosen Old English form \\
full \texttt{-i} stage retained too long & \texttt{*jugunþi} & expected
over-umlauted \texttt{y}-type result & negative control showing why
early \texttt{-i} loss must precede the OE umlaut stage \\
\end{longtable}
}

\protect\phantomsection\label{refs}
\begin{CSLReferences}{1}{1}
\bibitem[\citeproctext]{ref-BosworthToller1898}
Bosworth, Joseph, and T. Northcote Toller. 1898. \emph{An {Anglo-Saxon}
Dictionary}. Clarendon Press.

\bibitem[\citeproctext]{ref-BrightCassidyRingler1971}
Bright, James W. 1971. \emph{Bright's {Old English} Grammar and Reader}.
3rd edn. Edited by Frederic G. Cassidy and Richard N. Ringler. Holt,
Rinehart; Winston.

\bibitem[\citeproctext]{ref-Campbell1959}
Campbell, Alistair. 1959. \emph{Old {English} Grammar}. Clarendon Press.

\bibitem[\citeproctext]{ref-ClarkHall1960}
Clark Hall, J. R. 1960. \emph{A Concise {Anglo-Saxon} Dictionary}. 4th
edn. University of Toronto Press.

\bibitem[\citeproctext]{ref-Fulk2018}
Fulk, R. D. 2018. \emph{A Comparative Grammar of the Early {Germanic}
Languages}. Vol. 3. Studies in Germanic Linguistics. John Benjamins.

\bibitem[\citeproctext]{ref-KlugeSeebold2011}
Kluge, Friedrich. 2011. \emph{Etymologisches w{ö}rterbuch Der Deutschen
Sprache}. 25th edn. Edited by Elmar Seebold. De Gruyter.

\bibitem[\citeproctext]{ref-Kroonen2013}
Kroonen, Guus. 2013. \emph{Etymological Dictionary of {Proto-Germanic}}.
Vol. 11. Leiden Indo-European Etymological Dictionary Series. Brill.

\bibitem[\citeproctext]{ref-Luick1914}
Luick, Karl. 1914-\/-1940. \emph{Historische Grammatik Der Englischen
Sprache}. Tauchnitz.

\bibitem[\citeproctext]{ref-Orel2003}
Orel, Vladimir. 2003. \emph{A Handbook of {Germanic} Etymology}. Brill.

\bibitem[\citeproctext]{ref-RingeTaylor2014}
Ringe, Don, and Ann Taylor. 2014. \emph{The Development of {Old
English}}. Vol. 2. A Linguistic History of English. Oxford University
Press.

\bibitem[\citeproctext]{ref-SieversBrunner1965}
Sievers, Eduard. 1965. \emph{Altenglische Grammatik}. 3rd edn. Edited by
Karl Brunner. Niemeyer.

\end{CSLReferences}

\end{document}
